% GENERATED FILE. Edit canonical lessons, then run npm run generate:books.
% canonical-source-hash: fnv1a64:95c65e2afaa02caa
% canonical-lessons: IT-C434-rubinetto, IT-C434-gocciolare, IT-C434-intasato, IT-C434-scarico, IT-C434-sturare

\chapter{Il bagno non funziona}
\label{ch:it-il-bagno-non-funziona}
\hlchaptermodality{434}

\begin{chapteropening}
\textbf{By the end of this chapter:} \emph{I can report problems in the bathroom: the tap, dripping, a blocked sink, the drain and getting it unblocked.}

Complete the last lesson of chapter 434: I can report problems in the bathroom: the tap, dripping, a blocked sink, the drain and getting it unblocked.
\end{chapteropening}

\section[il rubinetto]{il rubinetto — the tap, the faucet}

\label{lesson:IT-C434-rubinetto}



\begin{quote}
Before the new one: say the Italian for the bar of soap, then the Italian for the hair dryer.
\end{quote}

\subsection*{You'll want to know: il rubinetto}
\textbf{il rubinetto} — \textquotedblleft{}the tap, the faucet\textquotedblright{}.

\begin{quote}
Il rubinetto dell'acqua calda non funziona.
\end{quote}

\begin{quote}
\emph{The hot tap does not work.}
\end{quote}

\emph{Aprire il rubinetto} — to turn on the tap.

\begin{cousinweb}
\emph{rubinetto} was borrowed from French \emph{robinet} 'tap' (EtymDB).
\end{cousinweb}

\begin{grammarlens}[title={What you've built}]
One more word for telling, explaining and describing at B1.
\end{grammarlens}

\subsection*{Your turn}
\begin{itemize}
\raggedright
  \item \emph{Say it:} \emph{il rubinetto}
  \item \emph{Say it:} \emph{il rubinetto}, once more
  \item \emph{Say it:} a sentence of your own with \emph{il rubinetto}
  \item \emph{Recall:} say the Italian for the bar of soap, then the Italian for the hair dryer, then say \emph{il rubinetto} again
\end{itemize}

\begin{tcolorbox}[breakable,colback=teal!4,colframe=teal!35!black,title={Before you move on}]
What does il rubinetto mean? (\textquotedblleft{}The tap, the faucet\textquotedblright{}.) Say it once more.
\end{tcolorbox}
\section[gocciolare]{gocciolare — to drip}

\label{lesson:IT-C434-gocciolare}



\begin{quote}
Before the new one: say the Italian for the hair dryer, then the Italian for the tap.
\end{quote}

\subsection*{You'll want to know: gocciolare}
\textbf{gocciolare} — \textquotedblleft{}to drip\textquotedblright{}.

\begin{quote}
Il rubinetto gocciola tutta la notte.
\end{quote}

\begin{quote}
\emph{The tap drips all night.}
\end{quote}

\emph{La goccia} — a drop.

\begin{cousinweb}
\emph{gocciolare} is built on \emph{goccia} 'drop', which is inherited from Latin \emph{gutta} 'drop' (Lewis \& Short).
\end{cousinweb}

\begin{grammarlens}[title={What you've built}]
One more word for telling, explaining and describing at B1.
\end{grammarlens}

\subsection*{Your turn}
\begin{itemize}
\raggedright
  \item \emph{Say it:} \emph{gocciolare}
  \item \emph{Say it:} \emph{gocciolare}, once more
  \item \emph{Say it:} a sentence of your own with \emph{gocciolare}
  \item \emph{Recall:} say the Italian for the hair dryer, then the Italian for the tap, then say \emph{gocciolare} again
\end{itemize}

\begin{tcolorbox}[breakable,colback=teal!4,colframe=teal!35!black,title={Before you move on}]
What does gocciolare mean? (\textquotedblleft{}To drip\textquotedblright{}.) Say it once more.
\end{tcolorbox}
\section[intasato]{intasato — blocked, clogged}

\label{lesson:IT-C434-intasato}



\begin{quote}
Before the new one: say the Italian for the tap, then the Italian for to drip.
\end{quote}

\subsection*{You'll want to know: intasato}
\textbf{intasato} — \textquotedblleft{}blocked, clogged\textquotedblright{}.

\begin{quote}
Il tubo è intasato.
\end{quote}

\begin{quote}
\emph{The pipe is blocked.}
\end{quote}

Roads too: \emph{le strade sono intasate}.

\begin{cousinweb}
\emph{intasato} is the past form of \emph{intasare} 'to clog', built on \emph{taso}, the crust that wine leaves in a barrel. EtymDB links \emph{taso} to the name of a Greek island known for its wine, but the origin is not settled.
\end{cousinweb}

\begin{grammarlens}[title={What you've built}]
One more word for telling, explaining and describing at B1.
\end{grammarlens}

\subsection*{Your turn}
\begin{itemize}
\raggedright
  \item \emph{Say it:} \emph{intasato}
  \item \emph{Say it:} \emph{intasato}, once more
  \item \emph{Say it:} a sentence of your own with \emph{intasato}
  \item \emph{Recall:} say the Italian for the tap, then the Italian for to drip, then say \emph{intasato} again
\end{itemize}

\begin{tcolorbox}[breakable,colback=teal!4,colframe=teal!35!black,title={Before you move on}]
What does intasato mean? (\textquotedblleft{}Blocked, clogged\textquotedblright{}.) Say it once more.
\end{tcolorbox}
\section[lo scarico]{lo scarico — the drain, the plughole; the flush}

\label{lesson:IT-C434-scarico}



\begin{quote}
Before the new one: say the Italian for to drip, then the Italian for blocked.
\end{quote}

\subsection*{You'll want to know: lo scarico}
\textbf{lo scarico} — \textquotedblleft{}the drain, the plughole; the flush\textquotedblright{}.

\begin{quote}
L'acqua non va giù nello scarico.
\end{quote}

\begin{quote}
\emph{The water will not go down the drain.}
\end{quote}

\emph{Il tubo di scarico} — the exhaust pipe of a car.

\begin{cousinweb}
\emph{scarico} is built on \emph{scaricare} 'to unload', from \emph{s-} and \emph{caricare} 'to load', which is inherited from Late Latin \emph{carricāre}, built on \emph{carrus} 'cart'.
\end{cousinweb}

\begin{grammarlens}[title={What you've built}]
One more word for telling, explaining and describing at B1.
\end{grammarlens}

\subsection*{Your turn}
\begin{itemize}
\raggedright
  \item \emph{Say it:} \emph{lo scarico}
  \item \emph{Say it:} \emph{lo scarico}, once more
  \item \emph{Say it:} a sentence of your own with \emph{lo scarico}
  \item \emph{Recall:} say the Italian for to drip, then the Italian for blocked, then say \emph{lo scarico} again
\end{itemize}

\begin{tcolorbox}[breakable,colback=teal!4,colframe=teal!35!black,title={Before you move on}]
What does lo scarico mean? (\textquotedblleft{}The drain, the plughole; the flush\textquotedblright{}.) Say it once more.
\end{tcolorbox}
\section[sturare]{sturare — to unblock, to unclog}

\label{lesson:IT-C434-sturare}



\begin{quote}
Before the new one: say the Italian for blocked, then the Italian for the drain.
\end{quote}

\subsection*{You'll want to know: sturare}
\textbf{sturare} — \textquotedblleft{}to unblock, to unclog\textquotedblright{}.

\begin{quote}
L'idraulico ha sturato lo scarico.
\end{quote}

\begin{quote}
\emph{The plumber unblocked the drain.}
\end{quote}

Also to uncork a bottle.

\begin{cousinweb}
\emph{sturare} is built from \emph{s-} and \emph{turare} 'to stop up', which is usually traced to Latin \emph{obtūrāre} 'to stop up' (Lewis \& Short), an inherited word.
\end{cousinweb}

\begin{grammarlens}[title={What you've built}]
One more word for telling, explaining and describing at B1.
\end{grammarlens}

\subsection*{Your turn}
\begin{itemize}
\raggedright
  \item \emph{Say it:} \emph{sturare}
  \item \emph{Say it:} \emph{sturare}, once more
  \item \emph{Say it:} a sentence of your own with \emph{sturare}
  \item \emph{Recall:} say the Italian for blocked, then the Italian for the drain, then say \emph{sturare} again
\end{itemize}

\begin{tcolorbox}[breakable,colback=teal!4,colframe=teal!35!black,title={Before you move on}]
What does sturare mean? (\textquotedblleft{}To unblock, to unclog\textquotedblright{}.) Say it once more.
\end{tcolorbox}
